\chapter{Qué hay que construir para decirlo hoy}
\label{cap:construir}

El argumento del capítulo anterior cabe en diez páginas de prosa. Formalizarlo —hacer que una
máquina lo verifique paso a paso— exige construir cinco cosas, ninguna de las cuales se puede
saltar. Este capítulo dice cuáles son y en qué orden vienen. Es, de paso, el mapa del libro.

\section{Las cinco piezas}

\begin{enumerate}[leftmargin=1.6em]
\item \textbf{Un lenguaje.} Qué es un término, qué es una fórmula, cómo se sustituye una variable.
  Suena trivial y no lo es: la sustitución con nombres exige evitar capturas, y ahí se pierden más
  demostraciones de las que uno espera.
\item \textbf{Una teoría.} Un lenguaje no afirma nada. Hacen falta axiomas — y decidir cuáles, que
  es una decisión con consecuencias que este libro persigue hasta el final.
\item \textbf{Un cálculo.} Qué cuenta como demostración. Y aquí aparece la primera sorpresa del
  proyecto: hacen falta \emph{dos}, por una razón que el capítulo siguiente explica.
\item \textbf{Una codificación.} Convertir fórmulas en números, de modo que la teoría pueda hablar
  de ellas.
\item \textbf{Un verificador.} Una fórmula del propio lenguaje que diga «esto de aquí es una
  demostración correcta». Sin ella no hay predicado de demostrabilidad, y sin predicado de
  demostrabilidad no hay teorema de Gödel.
\end{enumerate}

\begin{leccion}[El orden no es negociable]
Cada pieza usa las anteriores. No se puede codificar sin lenguaje, ni verificar sin codificación.
Por eso el libro va en ese orden, y por eso los capítulos técnicos empiezan diciendo qué se sabe ya
— si el lector llegó hasta aquí, tiene derecho a que se le diga dónde está.
\end{leccion}

\section{Lo que este libro añade a esa lista}

Si el libro fuera sólo la construcción, sería un manual más. Hay dos cosas que lo separan de un
manual, y conviene anunciarlas ahora porque condicionan cómo está escrito.

\textbf{La primera: aquí todo está compilado.} Ningún bloque de código de este libro se ha escrito a
mano. Se extraen del repositorio, y hay controles que fallan si el texto impreso deja de coincidir
con el código, si un teorema deja de estar demostrado desde la base declarada, o si el libro nombra
un símbolo que no existe. Debajo de cada enunciado hay una línea gris con los axiomas de los que
depende — el capítulo~\ref{cap:lean} explica qué es eso y por qué está ahí.

\textbf{La segunda: aquí también se cuentan los fracasos.} Una formalización real produce, además de
teoremas, una cantidad notable de callejones sin salida — y algunos de ellos enseñan más que los
teoremas. La Parte~\ref{parte:noescrito} está dedicada a eso: una inconsistencia que sobrevivió meses, cuatro
reparaciones que no funcionan, y la que sí. En los manuales eso no aparece, porque los manuales
presentan el resultado ya limpio.

\begin{leccion}[Por qué los fracasos son contenido y no anécdota]
Un resultado negativo demostrado —\emph{esto no se puede hacer, y aquí está la prueba}— ahorra al
siguiente el tiempo de intentarlo. En este proyecto hay varios, compilados. El libro los trata con
el mismo rango que los positivos, y tiene un entorno tipográfico propio para ellos.
\end{leccion}

\section{Cómo leer este libro}

Las partes de este libro siguen las cinco piezas, más los teoremas y más lo que no sale en los
libros:

\begin{center}\small
\begin{tabularx}{\linewidth}{@{}l >{\raggedright\arraybackslash}X@{}}
\toprule
\textbf{I} · El problema & dónde estás ahora: el contexto y las cuatro palabras de Lean que hacen falta \\
\textbf{II} · El terreno & el lenguaje, los cálculos y la teoría aritmética \\
\textbf{III} · Aritmetización & la codificación y el verificador \\
\textbf{IV} · Los teoremas & el punto fijo, Gödel~I, las condiciones de derivabilidad, Gödel~II \\
\textbf{V} · Lo que no sale en los libros & la inconsistencia, las reparaciones, el método \\
\bottomrule
\end{tabularx}
\end{center}

\noindent Los apéndices recogen el inventario de axiomas, el mapa de módulos, las trampas de Lean
encontradas, y un glosario y una tabla de notación a los que se puede volver en cualquier momento.
